%% ma: L12-B13-0021
%% lop: 12
%% chuong: 4
%% bai: 13
%% dang: 12.13.1
%% loai: TN4
%% muc: TH
%% dap_an: "C"
%% nguon: "(NBV)-ĐỀ SỐ 8-MỖI NGÀY 1 ĐỀ THI 2025.pdf (D08, MỖI NGÀY 1 ĐỀ - Đề số 8), Phần I Câu 7"
%% kien_thuc: "Bài 13 (thể tích khối tròn xoay quay quanh Ox), Bài 12 (tích phân)"
%% trang_thai: chinh-thuc
%% ghi_chu: "Đáp án gốc C đúng (sympy). Đề gốc nói 'trục hoành, trục tung và đường thẳng x = 2'; lời giải gốc dùng cận 0 đến 2 nên giữ nguyên. Gốc không có hình."
\begin{ex}
Hình phẳng $(H)$ giới hạn bởi đồ thị hàm số $y=\sqrt x$, trục hoành, trục tung và đường thẳng $x=2$. Vật thể tròn xoay được tạo ra khi quay $(H)$ quanh trục hoành có thể tích bằng
\choice{$\pi$}{$\dfrac{2\pi}3$}{\True $2\pi$}{$\dfrac{2\pi\sqrt2}3$}
\loigiai{$V=\pi\displaystyle\int_0^2\left(\sqrt x\right)^2\mathrm dx=\pi\int_0^2x\,\mathrm dx=\pi\cdot\dfrac{x^2}2\Big|_0^2=2\pi$.}
\end{ex}
